% ====================================================================
% §10 합산과 흑연 초기값 (N9)
% ====================================================================
\section{합산과 흑연 staging 초기값}\label{sec:sum}
전이별 peak 모양이 식~\eqref{eq:peakshape}(그 $L_V\to0$ 극한은 평형 종~\eqref{eq:tail-limit}\eqref{eq:eqpeak})로 정해지면, 용량 가중을 곱해 내부 전위 $V_n$ 위에 점별로 누적한다. 보존식 $Q_\cell q=Q_\bg+\sum_jQ_j\xi_j$ 를
$V$ 로 미분하면 서로소 자리 분해(전이 $j$ 의 진행이 서로 독립)에서 $\dd Q/\dd V=C_\bg+\sum_jQ_j\,\dd\xi_j/\dd V$ 이고,
각 $\dd\xi_j/\dd V$ 가 곧 peak 모양이라(식~\eqref{eq:peakshape}) 관측 ICA 는 배경과 전이별 기여의 선형 합으로 닫힌다:
\begin{equation}
\boxed{\;\frac{\dd Q}{\dd V}(V_n)\;=\;C_\bg\;+\;\sum_{j=1}^{N_p}Q_j\,
\big[\text{peak shape}_j\big]
\;=\;C_\bg+\sum_jQ_j\big[\text{평형 peak}-\text{지연 꼬리}\big]\;.}
\label{eq:sum}
\end{equation}
각 전이의 기여는 배경 위에 누적된다(배경은 합산 전에 먼저 깔림) --- 모든 평가가 입력 전위 $V_n$ 위에서 점별로 닫혀 출력 좌표가 곧 입력 좌표다. 스칼라 전위 입력은 한 점의 값으로 축약된다.

\subsection{staging 전이 초기값 --- 피팅 출발점}\label{sec:sum-staging}
표~\ref{tab:staging} 가 네 전이 각 인자의 의미와 출발값이다(신뢰값이 아닌 \emph{초기값} --- 사용자 피팅이
override 하는 것이 전제) --- 이 표가 있어 ``본 장만으로 곡선 재현 코드를 짤 수 있다''.

\begin{table}[t]
\centering\footnotesize
\renewcommand{\arraystretch}{1.3}
\caption{전이 초기값(흑연 staging 4건) --- 출발점, 피팅으로 override. $U$ 는 $(\Delta H_\rxn,\Delta S_\rxn)$
정합 목표. 전위 anchor: $2\!\to\!1$($0.085$ V, Dahn\cite{dahn1991})$\cdot$$2\mathrm
L\!\to\!2$($0.120$ V, Ohzuku 등\cite{ohzuku1993})는 보고값과 정확 일치(2차 출처 경유 확인이라 tier B 로
보수 분류), $4\!\to\!3$($0.210$ V, Dahn\cite{dahn1991})은 근사 확인(tier B),
$3\!\to\!2\mathrm L$($0.140$ V)는 Ohzuku 보고 $\approx0.125$ V\cite{ohzuku1993} 에서 $15$ mV 떨어진
배치값이다 --- 본 표의 두 anchor 출처\cite{dahn1991,ohzuku1993} 사이에서도 같은 전이의 보고 전위가
시료$\cdot$측정 조건에 따라 어긋나므로 그 시료 의존 편차 범위 안의 출발점으로 두고 tier C 로 분류한다. $\Delta S_\rxn$ 의 부호$\cdot$수십
J/(mol\,K) 스케일은 흑연 삽입 엔트로피 실측\cite{reynier2003} 과 정합(전이별 정밀값 아님 --- tier B).}
\label{tab:staging}
\begin{tabular}{l c c c c c c c}
\toprule
stage & $U$ [V] & $\Delta H_\rxn$ & $\Delta S_\rxn$ & $Q$ & $\Omega$ & $\Delta H_a$ & $|\dd V/\dd q|_{q_a}$ \\
 &  & [J/mol] & [J/(mol\,K)] & [$Q_\cell$] & [J/mol] & [J/mol] & [V] \\
\midrule
$4\!\to\!3$   & 0.210 & $-11700$ & $+29.0$ & 0.10 & 6000  & 48000 & 0.30 \\
$3\!\to\!2\mathrm L$ & 0.140 & $-13500$ & $0.0$   & 0.12 & 8000  & 46000 & 0.30 \\
$2\mathrm L\!\to\!2$ & 0.120 & $-13100$ & $-5.0$  & 0.25 & 10000 & 44000 & 0.30 \\
$2\!\to\!1$   & 0.085 & $-13000$ & $-16.0$ & 0.50 & 13000 & 40000 & 0.30 \\
\bottomrule
\end{tabular}
\end{table}

각 전이는 다중도 $n=1$, 폭 $w$ 폴백(0.020/0.016/0.014/0.012 V), $\Delta S_a=0$ 을 함께 갖는데, 폭
서식은 $n$ 이 \emph{우선} 적용되어 실제 초기 폭은 네 전이 균일하게 $w=RT/F\approx25.7$ mV 이며, $w$
폴백 열의 좁은 값은 전이 dict 에서 $n$ 입력을 제거할 때 활성화되는 대안 폭이다(\S\ref{sec:broadening}
② 의 같은 규정). $\Delta H_\rxn/\Delta S_\rxn$
은 $U(298)$ 정합으로 정해졌고($U=(-\Delta H_\rxn+298.15\Delta S_\rxn)/F$ 가 표의 $U$ 와 $\pm1$ mV 안에서 정합 --- $4\!\to\!3$ 행만 $0.2109$ 로 표시 자리수 경계), $\Delta H_a$ 는 저SOC(state of charge)$\to$
만충에서 감소하는 DFT(density functional theory) 활성화 경향\cite{persson2010,persson2010b} 에 맞춘
스케일이고(전이별 수치의 전용 문헌 anchor 는 근거 미발견 --- 실측 계면 활성화 에너지의 문헌 범위
$25$--$59$ kJ/mol 와 겹치는 출발점), $|\dd V/\dd q|_{q_a}$ 는 컷 OCV 기울기(피팅), $\Omega$ 는 정규용액
추정(전이별 문헌 anchor 없음 --- 피팅 출발점)이다.
초기 상태의 실현 크기를 수치로 확인하면 --- 표의 $\Delta H_a$ 초기값에서 대표 조건(298 K$\cdot$C/10)의
$L_{V,j}$ 는 $10^{-10}$--$10^{-8}$ V 규모로 폭 $w_j$(수십 mV)에 비해 무시할 만큼 작아 식~\eqref{eq:tail-limit} 의 $L_V\to0$ 극한이 사실상 이미 실현된다(본 수치는 $|I|/Q_\cell$ 를 $\mathrm s^{-1}$ 수치로 대입한 값 --- 정규화 $Q_\cell{=}1$ 에서 c-rate 숫자를 그대로 쓰는 규약; [1/h] 수치를 SI 로 환산하면 $\sim3600$ 배 작아지나 결론 $L_V\!\ll\!w$ 는 불변). 따라서 $\gamma$ 미배정($=0$)$\cdot$$R_n=0$ 과 함께 \emph{초기 상태의 곡선은 꼬리$\cdot$분기 없는 평형
종 전용}이다(가시적 꼬리는 피팅 개방 후 활성). $k_0=k_BT/h$$\cdot$$\Delta S_a{=}0$ 전제에서 꼬리가
폭 $w_j$ 에 견줄 만큼(가시적 꼬리) 자라는 데는 $\Delta H_a\approx80$ kJ/mol 급 또는 음의 $\Delta S_a$ 배정이 필요하므로
(전이별 --- 위 대표 조건과 $\chi=0.5$ 기본 기준), 인용 문헌
범위($25$--$59$ kJ/mol)는 절대값 anchor 가 아니라 경향 anchor 로 읽는다.
이 초기 상태로 합산식~\eqref{eq:sum} 을 실제로 평가한 전체 곡선 --- 문서가 최종 산출하는 바로 그 한 곡선 ---
이 그림~\ref{fig:sumcurve} 다: 배경 위 네 봉우리의 겹침과, 온도를 바꿀 때 봉우리들이 서로 반대 방향으로
움직이는 것(그림~\ref{fig:UjT} 의 기울기 부호)까지 한 장에서 읽힌다.

\begin{figure}[t]
\centering
\begin{tikzpicture}[x=46cm,y=0.60cm]
% ---- axes ----
\draw[-{Latex[length=1.6mm]}] (0.025,0) -- (0.025,9.6)
  node[above,font=\scriptsize] {$\dd Q/\dd V$ [$Q_\cell$/V]};
\draw[-{Latex[length=1.6mm]}] (0.025,0) -- (0.290,0)
  node[below,font=\scriptsize] {$V_n$ [V]};
\foreach \xx in {0.05,0.10,0.15,0.20,0.25}
  {\draw (\xx,0.10) -- (\xx,-0.10) node[below,font=\scriptsize] {\xx};}
\foreach \yy in {2,4,6,8}
  {\draw (0.0235,\yy) -- (0.0265,\yy) node[left,font=\scriptsize] {\yy};}
% ---- background C_bg ----
\draw[densely dotted] (0.028,0.05) -- (0.285,0.05);
\node[font=\tiny,anchor=west] at (0.031,0.16) {$C_\bg{=}0.05$};
% ---- per-transition contributions at 298.15 K (thin gray, on top of C_bg=0 baseline) ----
\draw[black!40] plot[smooth] coordinates {(0.03,1.815) (0.035,2.11) (0.04,2.432) (0.045,2.777) (0.05,3.137) (0.055,3.501) (0.06,3.854) (0.065,4.179) (0.07,4.458) (0.075,4.675) (0.08,4.814) (0.085,4.865) (0.09,4.825) (0.095,4.696) (0.1,4.488) (0.105,4.214) (0.11,3.894) (0.115,3.543) (0.12,3.18) (0.125,2.819) (0.13,2.472) (0.135,2.147) (0.14,1.849) (0.145,1.581) (0.15,1.343) (0.155,1.135) (0.16,0.9554)};
\draw[black!40] plot[smooth] coordinates {(0.05,0.5558) (0.055,0.6579) (0.06,0.7748) (0.065,0.907) (0.07,1.054) (0.075,1.215) (0.08,1.388) (0.085,1.568) (0.09,1.749) (0.095,1.926) (0.1,2.089) (0.105,2.229) (0.11,2.337) (0.115,2.407) (0.12,2.433) (0.125,2.413) (0.13,2.348) (0.135,2.244) (0.14,2.108) (0.145,1.948) (0.15,1.773) (0.155,1.591) (0.16,1.411) (0.165,1.237) (0.17,1.074) (0.175,0.9251) (0.18,0.791) (0.185,0.6721) (0.19,0.5682) (0.195,0.4782)};
\draw[black!40] plot[smooth] coordinates {(0.065,0.2276) (0.07,0.2705) (0.075,0.3201) (0.08,0.3768) (0.085,0.4408) (0.09,0.512) (0.095,0.5898) (0.1,0.673) (0.105,0.7595) (0.11,0.8467) (0.115,0.931) (0.12,1.008) (0.125,1.075) (0.13,1.125) (0.135,1.157) (0.14,1.168) (0.145,1.156) (0.15,1.124) (0.155,1.073) (0.16,1.006) (0.165,0.9283) (0.17,0.8439) (0.175,0.7566) (0.18,0.6702) (0.185,0.5872) (0.19,0.5096) (0.195,0.4386) (0.2,0.3748) (0.205,0.3183) (0.21,0.269)};
\draw[black!40] plot[smooth] coordinates {(0.14,0.2182) (0.145,0.2584) (0.15,0.3045) (0.155,0.3566) (0.16,0.4149) (0.165,0.4787) (0.17,0.5473) (0.175,0.619) (0.18,0.6918) (0.185,0.7627) (0.19,0.8286) (0.195,0.8858) (0.2,0.9307) (0.205,0.9604) (0.21,0.9728) (0.215,0.9668) (0.22,0.943) (0.225,0.9031) (0.23,0.8498) (0.235,0.7866) (0.24,0.717) (0.245,0.6445) (0.25,0.5721) (0.255,0.5022) (0.26,0.4366) (0.265,0.3763) (0.27,0.322) (0.275,0.2738) (0.28,0.2316)};
\node[black!40,font=\tiny] at (0.0715,5.05) {$2\!\to\!1$};
\node[black!40,font=\tiny] at (0.113,2.72) {$2\mathrm L\!\to\!2$};
\node[black!40,font=\tiny] at (0.145,1.47) {$3\!\to\!2\mathrm L$};
\node[black!40,font=\tiny] at (0.222,1.24) {$4\!\to\!3$};
% ---- total at 268.15 K (thick solid) ----
\draw[very thick] plot[smooth] coordinates {(0.03,1.675) (0.035,2.009) (0.04,2.398) (0.045,2.843) (0.05,3.343) (0.055,3.894) (0.06,4.485) (0.065,5.101) (0.07,5.72) (0.075,6.317) (0.08,6.865) (0.085,7.337) (0.09,7.712) (0.095,7.976) (0.1,8.122) (0.105,8.152) (0.11,8.071) (0.115,7.89) (0.12,7.622) (0.125,7.28) (0.13,6.878) (0.135,6.431) (0.14,5.955) (0.145,5.467) (0.15,4.981) (0.155,4.514) (0.16,4.078) (0.165,3.681) (0.17,3.33) (0.175,3.024) (0.18,2.76) (0.185,2.534) (0.19,2.336) (0.195,2.158) (0.2,1.993) (0.205,1.833) (0.21,1.676) (0.215,1.518) (0.22,1.362) (0.225,1.21) (0.23,1.063) (0.235,0.9254) (0.24,0.7987) (0.245,0.6844) (0.25,0.5831) (0.255,0.4948) (0.26,0.4189) (0.265,0.3544) (0.27,0.3) (0.275,0.2547) (0.28,0.217)};
% ---- total at 328.15 K (thick densely dashed) ----
\draw[very thick,densely dashed] plot[smooth] coordinates {(0.03,2.674) (0.035,3.034) (0.04,3.419) (0.045,3.823) (0.05,4.237) (0.055,4.652) (0.06,5.056) (0.065,5.437) (0.07,5.784) (0.075,6.084) (0.08,6.329) (0.085,6.513) (0.09,6.632) (0.095,6.686) (0.1,6.677) (0.105,6.609) (0.11,6.485) (0.115,6.313) (0.12,6.097) (0.125,5.843) (0.13,5.559) (0.135,5.252) (0.14,4.928) (0.145,4.596) (0.15,4.262) (0.155,3.935) (0.16,3.621) (0.165,3.325) (0.17,3.051) (0.175,2.803) (0.18,2.581) (0.185,2.386) (0.19,2.215) (0.195,2.066) (0.2,1.935) (0.205,1.818) (0.21,1.712) (0.215,1.612) (0.22,1.515) (0.225,1.419) (0.23,1.323) (0.235,1.225) (0.24,1.128) (0.245,1.03) (0.25,0.9346) (0.255,0.842) (0.26,0.7537) (0.265,0.6709) (0.27,0.5942) (0.275,0.5241) (0.28,0.4609)};
% ---- center-shift arrows (268 -> 328 K) ----
\draw[-{Latex[length=1.5mm]},thick] (0.0903,8.75) -- (0.0803,8.75);
\node[font=\tiny,anchor=south,align=center] at (0.0853,8.85)
  {$2\!\to\!1$: $-9.9$ mV\\($\Delta S_\rxn{<}0$)};
\draw[-{Latex[length=1.5mm]},thick] (0.2019,3.4) -- (0.2199,3.4);
\node[font=\tiny,anchor=south,align=center] at (0.211,3.5)
  {$4\!\to\!3$: $+18.0$ mV\\($\Delta S_\rxn{>}0$)};
\node[font=\tiny,anchor=west] at (0.228,3.0) {(중심 이동, $268\!\to\!328$ K)};
% ---- legend ----
\draw[very thick] (0.208,7.9) -- (0.220,7.9);
\node[font=\tiny,anchor=west] at (0.222,7.9) {총합, $268.15$ K};
\draw[very thick,densely dashed] (0.208,7.2) -- (0.220,7.2);
\node[font=\tiny,anchor=west] at (0.222,7.2) {총합, $328.15$ K};
\draw[black!40] (0.208,6.5) -- (0.220,6.5);
\node[font=\tiny,anchor=west] at (0.222,6.5) {전이별 기여, $298.15$ K};
\draw[densely dotted] (0.208,5.8) -- (0.220,5.8);
\node[font=\tiny,anchor=west] at (0.222,5.8) {배경 $C_\bg$};
\end{tikzpicture}
\caption{합산식~\eqref{eq:sum} 이 내는 전체 곡선 --- 좌표는 식 그대로의 수치 평가다: 초기 상태의 평형
종 전용 극한($L_V\!\to\!0$$\cdot$$\gamma_j{=}0$$\cdot$$R_n{=}0$)에서
$\dd Q/\dd V=C_\bg+\sum_jQ_j\,\xi_{\eq,j}(1-\xi_{\eq,j})/w_j$ 를 표~\ref{tab:staging} 초기값으로
평가했다($U_j(T)=(-\Delta H_{\rxn,j}+T\Delta S_{\rxn,j})/F$ 를 온도마다 재평가, $w_j{=}RT/F$
--- $n_j{=}1$ 균일, 배경은 상수 예시 $C_\bg{=}0.05\,Q_\cell$/V). 회색 가는 곡선이 $298.15$ K 의
전이별 기여, 굵은 실선$\cdot$파선이 $268.15/328.15$ K 의 총합이다. 온도가 오르면 (i) 폭
$w{=}RT/F$ 가 비례해 열려 봉우리가 낮고 넓어지고(면적 $Q_j$ 보존), (ii) 중심은
$\partial U_j/\partial T=\Delta S_{\rxn,j}/F$ 를 따라 \emph{전이마다 반대 방향}으로 움직인다 ---
$60$ K 동안 $2\!\to\!1$($\Delta S_\rxn{=}{-}16$ J/(mol\,K))은 $-9.9$ mV 로 왼쪽,
$4\!\to\!3$($+29$)은 $+18.0$ mV 로 오른쪽(화살표), $3\!\to\!2\mathrm L$($0$)은 제자리다. 다온도
피팅에서 온도마다 $U_j(T)$ 를 따로 읽어야 하는 까닭(\S\ref{sec:center})이 이 한 장이다.}
\label{fig:sumcurve}
\end{figure}
% 자산(v1.0.21 신규): [V21-R-03 합산 전체 곡선(도착 이미지) fig:sumcurve — FIGS_PICK N-5(베이스 FF3-3)]

\subsection{끝-대-끝 계산 예제 --- $V_n=0.085$ V 한 점}\label{sec:sum-worked}
``본 장만 보고 같은 곡선을 재현한다''는 선언을 한 점에서 실증한다. 설정은 초기 상태 그대로다:
$298.15$ K, 평형 종 전용($L_V\!\to\!0$·$\gamma_j{=}0$·$R_n{=}0$), 표~\ref{tab:staging} 입력,
배경 $C_\bg=0.05\,Q_\cell$/V(그림~\ref{fig:sumcurve} 와 같은 예시 상수), 평가점 $V_n=0.085$ V.

\textbf{(a) 폭과 중심.} $n_j{=}1$ 균일이라 폭은 네 전이 공통 $w=RT/F=25.693$ mV. 중심은
식~\eqref{eq:Uj} 에 표의 $(\Delta H_\rxn,\Delta S_\rxn)$ 를 넣은 환산값을 쓴다 --- 표시 반올림 $U$ 열
($0.210/0.140/0.120/0.085$ V)을 되넣지 않는다(반올림 입력은 아래 (b) 의 지수에서 mV 급 오차를 낳는다):
$U_j(298.15)=210.88/139.92/120.32/85.30$ mV.

\textbf{(b) 전이별 평형 진행률.} 식~\eqref{eq:logisticsolve} 의 무차원 인자 $(V_n-U_j)/w$ 와 진행률:
\[
\begin{array}{lcccc}
 & 4\!\to\!3 & 3\!\to\!2\mathrm L & 2\mathrm L\!\to\!2 & 2\!\to\!1\\
(V_n-U_j)/w & -4.899 & -2.137 & -1.375 & -0.011\\
\xi_{\eq,j} & 0.0074 & 0.1055 & 0.2019 & 0.4971
\end{array}
\]
--- 평가점이 stage $2\!\to\!1$ 중심 바로 옆이라 그 전이만 반쯤 진행했고, 나머지는 초입이다.

\textbf{(c) 전이별 봉우리 몫.} 식~\eqref{eq:eqpeak} 의 $Q_j\,\xi_{\eq,j}(1-\xi_{\eq,j})/w_j$ [$Q_\cell$/V]:
\[
0.0286\;(4\!\to\!3)\quad 0.4408\;(3\!\to\!2\mathrm L)\quad 1.568\;(2\mathrm L\!\to\!2)\quad 4.865\;(2\!\to\!1).
\]

\textbf{(d) 합산.} 식~\eqref{eq:sum} 으로
\[
\frac{\dd Q}{\dd V}\Big|_{0.085\,\mathrm V}=C_\bg+\sum_j Q_j\frac{\xi_{\eq,j}(1-\xi_{\eq,j})}{w_j}
=0.05+6.902=6.95\ Q_\cell/\mathrm{V}.
\]
\begin{verifybox}
검산 --- (i) $2\!\to\!1$ 몫: $\xi_\eq\simeq\tfrac12$(중심)이라 순높이가 식~\eqref{eq:eqpeak} 의 읽기값
$Q_j/(4w_j)=0.50/(4\times0.025693)=4.87$ 과 일치한다. (ii) 그림 대조: (c) 중 도면에 그려진 세 값($3\!\to\!2\mathrm L\cdot2\mathrm L\!\to\!2\cdot2\!\to\!1$)은
그림~\ref{fig:sumcurve} 의 $298.15$ K 전이별 기여 곡선을 $V_n=0.085$ V 에서 읽은 값 그대로이고, $4\!\to\!3$ 몫 $0.0286$ 은 그 곡선의 표시역($0.14$ V 이상) 밖이라 도면에는 나타나지 않는다(수치 대조만 --- 독립
수치 평가와 교차 일치). (iii) 재현 절차: 본 예제에 쓰인 것은 표~\ref{tab:staging} 와
식~\eqref{eq:Uj}·\eqref{eq:logisticsolve}·\eqref{eq:eqpeak}·\eqref{eq:sum} 뿐 --- 다른 입력이 없으므로,
같은 네 식을 전 $V_n$ 축에 돌리면 그림~\ref{fig:sumcurve} 전체가 재현된다. 가역 발열까지 잇는 같은
정신의 예제는 Part T 의 계산 예제(\S\ref{ssec:worked})가 담당한다.
\end{verifybox}

% 자산(v1.0.21 신규): [V21-Q5-01 끝-대-끝 계산 예제 sec:sum-worked — D21-2 top3 ①(재현 선언 실증·B-006 환산값 규약 재강조)]

여기까지가 흑연 곡선(Part I)의 결과 사슬이다 --- 이어지는 Part T 는 같은 전극$\cdot$같은 입력에서 곡선
대신 열(엔트로피 분해와 가역 발열)을 뽑고, 이 골격을 두 번째 전극 LCO 양극에 거는 일은 Chapter 2 가 맡는다.

% 자산: [B2-030] [B2-031] [B2-032] [B2-033] [B2-034]
